\section{Habitat Fragmentation Index}

\label{sec:fragmentation}

\subsection{Definition}

The HFI quantifies the degree to which detected land use changes degrade habitat connectivity within a defined landscape unit.
For a landscape $L$ at time $t$, the HFI is computed as:

\begin{equation}
    \text{HFI}(L, t) = \alpha \cdot \Delta\text{PC}(L, t) + \beta \cdot \Delta\text{ED}(L, t) + \gamma \cdot \Delta\text{COHESION}(L, t)
\end{equation}

\noindent where:
\begin{itemize}[leftmargin=*]
    \item $\Delta\text{PC}$ is the change in Probability of Connectivity \citep{saura2007new}, measuring functional connectivity loss.
    \item $\Delta\text{ED}$ is the change in Edge Density, measuring the increase in habitat-non-habitat interfaces.
    \item $\Delta\text{COHESION}$ is the change in Patch Cohesion Index, measuring the physical connectedness of habitat patches.
\end{itemize}

The weights $\alpha, \beta, \gamma$ are calibrated against field-measured biodiversity indicators (species richness, abundance indices) from 87 monitoring sites with concurrent satellite and ground data.

\subsection{Connectivity Graph}

The Probability of Connectivity (PC) is computed on a habitat connectivity graph $G = (V, E)$ where:
\begin{itemize}[leftmargin=*]
    \item Nodes $V$ represent habitat patches identified from the land cover classification.
    \item Edges $E$ connect patches within species-specific dispersal distances.
    \item Edge weights $w_{ij}$ represent dispersal probability, computed from a least-cost path analysis through the intervening land cover matrix.
\end{itemize}

\begin{equation}
    \text{PC} = \frac{\sum_{i=1}^{n} \sum_{j=1}^{n} a_i \cdot a_j \cdot p_{ij}^*}{A_L^2}
\end{equation}

\noindent where $a_i$ is the area of patch $i$, $p_{ij}^*$ is the maximum product probability of dispersal between patches $i$ and $j$, and $A_L$ is the total landscape area.

\subsection{Temporal HFI Trends}

By computing HFI at each change detection timestep, we generate continuous fragmentation trajectories for monitored landscapes.
The rate of change $\dot{\text{HFI}}$ provides an early warning indicator:

\begin{equation}
    \dot{\text{HFI}}(L, t) = \frac{d}{dt}\text{HFI}(L, t) \approx \frac{\text{HFI}(L, t) - \text{HFI}(L, t - \Delta t)}{\Delta t}
\end{equation}

An accelerating HFI ($\ddot{\text{HFI}} > 0$) indicates that fragmentation is intensifying, triggering heightened monitoring and alert priority.

\subsection{Validation Against Biodiversity Metrics}

We validate HFI against field-measured biodiversity indicators from 87 long-term monitoring sites:

\begin{table}[H]
\centering
\caption{Correlation between HFI and field biodiversity metrics.}
\label{tab:hfi_validation}
\begin{tabular}{@{}lcc@{}}
\toprule
\textbf{Biodiversity Metric} & \textbf{Pearson $r$} & \textbf{$p$-value} \\
\midrule
Species Richness Change & $-0.87$ & $< 0.001$ \\
Shannon Diversity Index Change & $-0.82$ & $< 0.001$ \\
Population Abundance Index & $-0.79$ & $< 0.001$ \\
Functional Diversity Change & $-0.74$ & $< 0.001$ \\
Genetic Diversity (microsatellites) & $-0.68$ & $< 0.001$ \\
\bottomrule
\end{tabular}
\end{table}

HFI correlates most strongly with species richness change ($r = -0.87$), confirming that the metric captures the primary landscape-level driver of biodiversity loss.
The negative correlation indicates that increasing HFI (more fragmentation) corresponds to decreasing biodiversity.

% ============================================================
% 5. Conservation Alert System
% ============================================================
